\documentclass[11pt,a4paper,twoside,notitlepage]{article}
%comment out the line above and decomment the one below; the format will change
%\documentclass[conference,a4paper]{IEEEtran}  

%
% Here, you specify the packages you need
%

%You need these 2 packages to write algorithms
\usepackage{algorithmic}
\usepackage{algorithm}
%You need this package to handle figures
\usepackage{graphics}
%You need this package to write Romanian diacritics properly
\usepackage{combelow}

\usepackage[utf8x]{inputenc} %alternative solution for Romanian diacritics
%This is how you define you own commands/macros
%this is a macro for writing "n choose k" the Romanian style
\newcommand{\mycomb}[2]{{C}_{#1}^{#2}} 
%this is a macro for writing "n choose k" the English style
%\newcommand{\mycomb}[2]{\left( \begin{array}{l}{#2}\\{#1}\\ \end{array} \right)} %this 


%
% The document content starts here
%

\begin{document}
% paper title 
\title{\LaTeX -- let's roll!}

% author names and affiliations
\author{
Radu R\u{a}zvan Sl\u{a}vescu\\ %double backslash means linebreak
Rom\^{a}nia\\
}

%generate paper title
\maketitle 

%this is the document's abstract
\begin{abstract}
We will present how to write a \texttt{.tex} file and how to generate a \texttt{.pdf} from it.
This includes writting text, equations, tables, algorithms and handling references.
You need little more than one hour to understand the basics of \LaTeX\ and one life to master it :)
\end{abstract} 

%This is how you start a section
\section{Introduction}

You should look in the \texttt{.tex} file (which is just a text file containing your text plus some formatting macros)
as well as in the \texttt{.pdf} file generated from it. 
You will see plain content and some macros. 
A macro always starts with a backslash. For example, the macro \verb|\section{...}| %this is one way of entering verbatim text
marks the beginning of a document section (of course, you should replace the dots with the section's actual name). 

The comments are introduced by the \% sign and will help you understand the meaning of the macros; the comment continues up to the end of line. 

The commands needed for generating the \texttt{.pdf} are written in the file called \texttt{Makefile}; 
make sure the variables in its first lines are properly set. 
Then, all you need to do is to type \texttt{make} at the shell prompt.
Alternatively, you can use \LaTeX\ %this is how we properly write the name of the magical LaTeX. 
editors (see below) which provide the compiling commands. 

\subsection{Some basics}  %This introduces a subsection. Latex will number it for you, e.g., 1.1

%We will present how to write text, equations, tables, algorithms and how to handle references.


First, Romanian letters:
\u{a} 
\^{\i}
\cb{s} 
\cb{t} 
\u{A} 
\^{I}
\cb{S} 
\cb{T} 
\^a 
\^A
\"{u} (even if the latter one does not seem to be in Romanian). 
In order to write them properly (use a comma instead of a cedilla for \cb{s}), you need to add \verb|\usepackage{combelow}| to your document's preamble. 
Or you can add \verb|\usepackage[utf8x]{inputenc}| to your document's preamble, do a \verb|setxkbmap ro| in a terminal to change the keyboard layout 
to Romanian, then you can enter Romanian diacritics using \texttt{AltGr} (e.g., \texttt{AltGr+t} for ț). 

We can use bullets:
\begin{itemize}
    \item number one
    \item number two
\end{itemize}

Or we can use numbered items: 
\begin{enumerate}
    \item un (aka 1)
    \item dos (aka 2)
\end{enumerate}

To start a new paragraph, you just enter a blank line in the \texttt{.tex} file. The first paragraph will have no indentation. 
This is normal; do not attempt to alter this.

Some characters, like \$, \&, \%, \#, \_, \{, \}, have special meanings in \LaTeX; if you want to escape them, place a backslash in front of them. 
In order to write a backslash in your \texttt{.pdf}, put \verb|\backslash| in your \texttt{.tex} file. To write $<$ or $>$, you should enclose them 
between \$ signs (\verb|$<$| and \verb|$>$| respectively). 

When we cite a paper (e.g., paper \cite{Zhou04} %this macro helps you cite
), we take its bibtex reference from the web,
store it in the \texttt{.bib} file and label it there (e.g., \texttt{gigel}; this label must be unique in the \texttt{.bib} file).
Then use \verb|\cite{gigel}| in the \texttt{.tex} file and let \LaTeX\ handle ordering and cross-referencing. 
Each time you use an idea, text etc. borrowed from some other author, you must properly cite the source. 
It is a good time now to take a look in the \texttt{.bib} file and see how a bibtex entry of an article looks like.
Then, you may want to search the web for the bibtex of an article (e.g., search bibtex "The Anatomy of a Large-Scale Hypertextual Web Search Engine"), 
grab it and add it to the \texttt{.bib} file.

The rest of this paper is structured as follows. 
Section \ref{sec:s22} %crossreferences; see below, where the label is introduced
describes the implementation details, etc. 


\section{Implementation Details} 
\label{sec:s22} %this is what you need to write when you reffer to this section; here you just introduced the label

If we have algorithms used/modified/implemented/introduced, we can describe them. 
For instance, algorithm \ref{ftalg}. 
Describing algorithms requires studying the documentation of the \texttt{algorithmic} and \texttt{algorithm} packages.

\begin{algorithm}
 \caption{An algorithm looks like this} %algorithm name
   \label{ftalg} %label used to reffer it from text
   \begin{algorithmic}
     \IF{$Committed(G_{1},GR,\alpha)$} 
      \STATE $BRT_{\alpha}$ = $PredictBRT(G_{1},GR,\alpha,C_{GR})$
      \STATE $C_{\beta}$ = $ContextUpdate(C_{\beta},o)$
     \ENDIF
      \IF{$utility \ge CommunicationCost(G_{2})$}
        \STATE $Int.To(G_{1},Communicate(G_{1},G_{2},o))$
      \ENDIF
   \end{algorithmic}
\end{algorithm}

We can write equations:
% this is an equation 
\begin{equation} 
  \frac{S:a_{m} \land r_{b} \land d}
       {B:d_{c} \land ((d_{e} \land \neg d_{t_{2}} \land d_{t_{15}}) \lor (\neg d_{e} \land d_{t_{2}})}
\label{eq:e50}
\end{equation}

And big curly brackets:

\[\left \{ 
    \begin{array}{l}
        b=\frac{r}{r+s+2} \\
        d = \frac{s}{r+s+2}\\ 
        u = \frac{2}{r+s+2}
    \end{array}\right.
\]

The notations for "n choose k" and sum look like 
%inside a text line, mathematical formulas are written between dollar signs $
$\mycomb{n}{k}$ %Using a macro. Change the implementation in the document's preamble and the formula layout will change - no modification inside the document
 and $\displaystyle \sum_{k=0}^{n} \mycomb{n}{k}$ respectively.

\section{Results}

\subsection{Tables}

We summarize the results in a table, e.g., Table \ref{Mom}. Typically, the table is placed on top of the page where it is reffered for the first time, 
or on one of the consecutive pages. It is the \LaTeX\  compiler decision. 

\begin{table}% macros for table handling
    \caption{Multinomial opinion multiplication}
    \label{Mom}
        \begin{center}
            \begin{tabular}{|l|l|l|l|l|l|l|}
                \hline 
                & \multicolumn{3}{c|}{belief} & \multicolumn{3}{c|}{atomicity} \\
                \cline{2-7}
                & $poor$ & $avg$ & $good$ & $poor$ & $avg$ & $good$ \\ 
        	\hline
                $success$   &  0.1    &  0.2    &  0.3    & 0.4    & 0.5    &  0.6     \\
                $failure$   &  0.1    &  0.2    &  0.3    & 0.4    & 0.5    &  0.6     \\
                                                    \hline
	    \end{tabular}
        \end{center}
\end{table}

\subsection{Figures}
We can also have figures, like Figure \ref{fig:bye}. Some sources claim they tell 1000 words/figure. 
 %this is how you add a picture 
 \begin{figure}
  \begin{center}
   \includegraphics{img/theend2.eps}
   \caption{Saying bye bye} %title figure
   \label{fig:bye} %figure label for reffering from text
  \end{center}
 \end{figure}
 


\section{Conclusions}

Now we know how to write text, tables, references etc.


\section*{Acknowledgments} % use section* for acknowledgement as you don't want to number it
Special thanks to Knuth and Lamport for section \ref{sec:s22}. %this is how you refer to a section


\appendix

\section{How to run the example}


The \texttt{v1.tex} file contains the source; edit it with
%another way of introducing verbatim text
\begin{verbatim} 
kile v1.tex
\end{verbatim}

To generate the \texttt{.pdf} file, type
\begin{verbatim} 
make
\end{verbatim}

And no, we don't attempt to launch the \texttt{Makefile}. Just type \texttt{make} at the shell prompt and enjoy. 

Alternatively, you may press the QuickBuild button in the Kile editor.  

%this is to generate the bibliography
\bibliographystyle{IEEEtran} %same as .bst
\bibliography{v1} %same file name as bib

% that's all folks
\end{document}
